\chapter{Datos y material de trabajo}

\section{Fuente de datos y datos obtenidos}

El sistema trabaja con datos financieros históricos descargados previamente mediante \texttt{yfinance}~\cite{yfinance} y almacenados en ficheros CSV. Yahoo Finance publica información de mercado, fichas de activos y series históricas para acciones, fondos cotizados, índices, divisas, criptomonedas y materias primas. La librería \texttt{yfinance} permite acceder de forma programática a parte de esta información disponible en Yahoo Finance~\cite{yahoofinancehistorical}, incluyendo campos habituales como apertura, máximo, mínimo, cierre, cierre ajustado y volumen. En este trabajo se utiliza como herramienta práctica de descarga para un contexto académico, no como proveedor garantizado para producción.

Cada CSV se conserva junto a un fichero de metadatos en formato JSON. Este metadato registra la consulta original, el objetivo analítico, los símbolos financieros implicados, el intervalo temporal y los parámetros de descarga. La decisión permite separar la fase de obtención de datos de la fase de análisis: el MVP no depende de que la API externa esté disponible durante cada ejecución, sino que opera sobre un conjunto de datos congelado y revisable.

La Tabla~\ref{tab:datos_obtenidos} muestra los datos obtenidos y utilizados como base del TFM. No se incluyen todos los nombres de fichero en la tabla para evitar una presentación demasiado técnica; los nombres completos se mantienen en \texttt{data/raw} y en el resumen procesado \texttt{data/processed/eda\_resumen\_ficheros.csv}.

\begin{table}[H]
\centering
\small
\begin{tabularx}{\textwidth}{L{0.22\textwidth}L{0.19\textwidth}L{0.16\textwidth}Y}
\toprule
\textbf{Caso de datos} & \textbf{Activos} & \textbf{Cobertura} & \textbf{Uso en el MVP} \\
\midrule
NVDA cinco años & NVDA & 2021--2026, diario & Crecimiento de precio y validación de respuestas de Nivel A. \\
NVDA frente a AMD & NVDA, AMD & 2024--2026, diario & Comparación entre activos y cálculo de rentabilidad relativa. \\
AAPL tres meses & AAPL & 2025--2026, diario & Visión descriptiva de un activo y análisis técnico básico. \\
QQQ frente a SPY & QQQ, SPY & 2024, diario & Retorno, riesgo, drawdown y comparación estructurada. \\
S\&P 500 desde 2020 & \texttt{\^{}GSPC} & 2020--2026, diario & Serie amplia de mercado para pruebas de histórico largo. \\
Bitcoin 2024 & BTC-USD & 2024, diario & Activo cripto con alta variabilidad. \\
EUR/USD horario & EURUSD=X & 10 días, 1h & Ejemplo de divisa con frecuencia intradía. \\
Oro horario & GC=F & 1 semana, 1h & Materia prima con frecuencia intradía. \\
\bottomrule
\end{tabularx}
\caption{Datos obtenidos y finalidad de cada conjunto dentro del MVP.}
\label{tab:datos_obtenidos}
\end{table}

La variedad de activos permite comprobar que el sistema no está limitado a una acción concreta. También fuerza al flujo LLM a trabajar con horizontes distintos, desde consultas diarias de varios años hasta series horarias de pocos días.

\section{Catálogo de consultas y proceso de descarga}

El punto de partida del conjunto de datos es un catálogo de consultas de referencia. Cada entrada define la consulta, los tickers, el rango temporal o período relativo, el intervalo de observación y los posibles avisos. Este catálogo actúa como contrato entre la fase de preparación de datos y el flujo de análisis.

El proceso utilizado para construir el conjunto de datos puede resumirse en cuatro pasos:

\begin{enumerate}
    \item Definir o revisar la consulta en el catálogo, por ejemplo una comparativa entre QQQ y SPY o una descarga de NVDA a cinco años.
    \item Transformar esa consulta en parámetros de descarga: lista de tickers, \texttt{start}, \texttt{end}, \texttt{period} e \texttt{interval}.
    \item Descargar las series con \texttt{yfinance} y guardar el resultado en \texttt{data/raw} como CSV, junto a un \texttt{.metadata.json} con los parámetros usados.
    \item Normalizar los CSV, generar tablas procesadas y producir resúmenes exploratorios auxiliares.
\end{enumerate}

En el repositorio actual, los CSV crudos ya están congelados. Por tanto, para reproducir la parte documentada de esta memoria no es necesario volver a llamar a Yahoo Finance: basta con conservar los datos crudos y regenerar las tablas procesadas si fuera necesario. Si se quisiera ampliar el conjunto de datos, el procedimiento correcto sería añadir primero la consulta al catálogo y después generar un nuevo CSV con su metadato asociado.

Cada consulta se transforma en una representación estructurada antes de entrar en el flujo LLM. El siguiente bloque muestra un ejemplo simplificado:

\begin{small}
\begin{verbatim}
{
  "query": "Compárame QQQ y SPY en 2024 de forma clara,
            con tabla comparativa y serie normalizada",
  "tickers": ["QQQ", "SPY"],
  "csv_paths": [
    "data/raw/qqq_y_spy_desde_20240101_hasta_20241231.csv"
  ],
  "start": "2024-01-01",
  "end": "2024-12-31",
  "period": null,
  "interval": "1d",
  "warnings": []
}
\end{verbatim}
\end{small}

La Tabla~\ref{tab:formato_consulta} resume el significado de los campos principales.

\begin{table}[H]
\centering
\small
\begin{tabularx}{\textwidth}{L{0.25\textwidth}Y L{0.23\textwidth}}
\toprule
\textbf{Elemento} & \textbf{Descripción} & \textbf{Ejemplo} \\
\midrule
Consulta original & Texto introducido o definido como caso de prueba. & Comparar Nvidia y AMD. \\
Objetivo analítico & Tipo de análisis que debe ejecutarse. & Comparación de activos. \\
Símbolos financieros & Activos sobre los que se realiza el análisis. & NVDA, AMD. \\
Rango temporal & Período relativo o fechas de inicio y fin. & Dos años. \\
Intervalo & Frecuencia de los datos históricos. & Diario. \\
Avisos & Observaciones detectadas al preparar la entrada. & Sin avisos. \\
\bottomrule
\end{tabularx}
\caption{Estructura semántica de una consulta financiera antes de su procesamiento.}
\label{tab:formato_consulta}
\end{table}

\section{Estructura de los CSV}

Los ficheros generados por \texttt{yfinance} contienen columnas de tipo OHLCV: apertura, máximo, mínimo, cierre, cierre ajustado y volumen. Cuando hay más de un ticker, las columnas se guardan con cabeceras multinivel: el primer nivel identifica el símbolo financiero y el segundo nivel identifica el campo de mercado. Por este motivo, los CSV no se leen como tablas planas con una sola fila de cabecera, sino con \texttt{header=[0, 1]}.

La Tabla~\ref{tab:estructura_csv} resume los campos usados por el sistema. La columna de cierre es la variable principal, porque de ella dependen las métricas de crecimiento, rentabilidad, volatilidad, drawdown e indicadores técnicos básicos.

\begin{table}[H]
\centering
\small
\begin{tabularx}{\textwidth}{L{0.18\textwidth}L{0.25\textwidth}Y}
\toprule
\textbf{Campo} & \textbf{Uso principal} & \textbf{Observación} \\
\midrule
\texttt{Open} & Precio de apertura. & Se conserva como contexto, aunque no es obligatorio en todas las respuestas. \\
\texttt{High} y \texttt{Low} & Rango máximo y mínimo. & Útiles para revisar amplitud diaria o intradía. \\
\texttt{Close} & Cálculo financiero central. & Base de rentabilidad, volatilidad, drawdown y medias móviles. \\
\texttt{Adj Close} & Precio ajustado. & Se conserva si aparece en la descarga original. \\
\texttt{Volume} & Actividad negociada. & Puede ser cero o menos informativo en divisas y algunos futuros. \\
\bottomrule
\end{tabularx}
\caption{Campos OHLCV presentes en los CSV de mercado.}
\label{tab:estructura_csv}
\end{table}

La estructura multinivel también explica por qué el proyecto incluye una capa de normalización. El proceso exploratorio transforma cada combinación caso--ticker en formato largo, con columnas explícitas para identificador de caso, fichero fuente, ticker, intervalo, fecha y campos OHLCV. Este formato facilita auditar métricas por serie y evita que el código generado por el LLM tenga que resolver por sí solo todas las variantes de cabecera.

\section{Análisis exploratorio de datos}

El análisis exploratorio lee los CSV descargados con \texttt{yfinance}, interpreta las cabeceras multinivel, normaliza los datos a formato largo y genera resúmenes procesados por fichero y por serie.

El análisis confirma que el conjunto contiene 5.080 filas normalizadas, con un rango observado entre el 2 de enero de 2020 y el 17 de marzo de 2026. Los intervalos considerados son diario (\texttt{1d}) e intradía horario (\texttt{1h}).

\begin{table}[H]
\centering
\small
\begin{tabular}{lr}
\toprule
\textbf{Indicador} & \textbf{Resultado} \\
\midrule
Casos definidos en catálogo & 10 \\
CSV analizados & 8 \\
Series por símbolo y caso & 10 \\
Filas normalizadas & 5.080 \\
Nulos en precios de cierre & 0 \\
Duplicados por fecha & 0 \\
\bottomrule
\end{tabular}
\caption{Resumen de cobertura y calidad del conjunto de datos.}
\label{tab:resumen_cobertura_datos}
\end{table}

\begin{table}[H]
\centering
\small
\begin{tabularx}{\textwidth}{L{0.28\textwidth}L{0.18\textwidth}rL{0.26\textwidth}}
\toprule
\textbf{Caso} & \textbf{Tickers} & \textbf{Filas} & \textbf{Rango observado} \\
\midrule
AAPL tres meses & AAPL & 60 & 2025-12-17 a 2026-03-16 \\
NVDA frente a AMD & AMD, NVDA & 500 & 2024-03-18 a 2026-03-16 \\
NVDA cinco años & NVDA & 1.255 & 2021-03-17 a 2026-03-16 \\
Bitcoin 2024 & BTC-USD & 365 & 2024-01-01 a 2024-12-30 \\
S\&P 500 desde 2020 & \texttt{\^{}GSPC} & 1.558 & 2020-01-02 a 2026-03-16 \\
QQQ frente a SPY & QQQ, SPY & 251 & 2024-01-02 a 2024-12-30 \\
\bottomrule
\end{tabularx}
\caption{Muestra de cobertura temporal obtenida en el EDA.}
\label{tab:cobertura_eda}
\end{table}

Además del resumen agregado, se calcularon métricas descriptivas por serie. La Tabla~\ref{tab:resumen_series_eda} muestra una selección de resultados. Se observan activos con crecimientos muy elevados, como NVDA en cinco años, series con caídas acumuladas significativas, como AMD en la comparativa con NVDA, y activos de menor volatilidad relativa, como SPY. Esta diversidad obliga al sistema a redactar respuestas para escenarios con magnitudes y niveles de riesgo distintos.

\begin{table}[H]
\centering
\small
\begin{tabularx}{\textwidth}{L{0.26\textwidth}rrr}
\toprule
\textbf{Serie} & \textbf{Retorno total} & \textbf{Volatilidad anual} & \textbf{Máx. caída} \\
\midrule
NVDA, 5 años & 1273.33 \% & 51.31 \% & -66.36 \% \\
BTC-USD, 2024 & 109.76 \% & 44.13 \% & -26.18 \% \\
NVDA, 2 años & 107.13 \% & 49.42 \% & -36.89 \% \\
S\&P 500, desde 2020 & 105.64 \% & 20.77 \% & -33.93 \% \\
QQQ, 2024 & 28.07 \% & 17.99 \% & -13.56 \% \\
SPY, 2024 & 24.45 \% & 12.63 \% & -8.41 \% \\
AMD, 2 años & 3.11 \% & 55.10 \% & -58.98 \% \\
AAPL, 3 meses & -7.00 \% & 24.15 \% & -10.07 \% \\
\bottomrule
\end{tabularx}
\caption{Métricas exploratorias por serie utilizadas para contextualizar los casos de prueba.}
\label{tab:resumen_series_eda}
\end{table}

\section{Calidad, cobertura y limitaciones}

La ausencia de nulos en el precio de cierre es relevante porque las métricas principales del MVP dependen directamente de esta columna. También se comprobó la ausencia de duplicados por fecha en los ficheros analizados. Estas condiciones permiten ejecutar las consultas actuales sin introducir imputaciones adicionales.

Como limitación, las descargas basadas en períodos relativos, por ejemplo tres meses o cinco años, dependen de la fecha en la que se generó el CSV. Para una versión final completamente cerrada sería recomendable fijar fechas explícitas de inicio y fin en todos los casos. También debe considerarse que los datos proceden de Yahoo Finance mediante \texttt{yfinance}. Al tratarse de una herramienta abierta no afiliada ni avalada por Yahoo, su uso se considera adecuado para un contexto académico y experimental, pero en un entorno productivo habría que estudiar licencias, disponibilidad, estabilidad del proveedor y derechos de uso de los datos descargados~\cite{yfinance}.

\section{Material generado por el flujo y trazabilidad}

Además de los CSV históricos, el proyecto genera material intermedio durante cada ejecución. Este material es importante porque permite revisar cómo una consulta en lenguaje natural se transforma en una respuesta final. En un TFM no basta con indicar que el sistema responde: también debe poder explicarse qué entendió el modelo, qué cálculo se preparó, qué datos se usaron y qué resultado recibió el agente que redacta la explicación.

La Tabla~\ref{tab:artefactos_flujo_agentes} resume los artefactos conceptuales del flujo. La nomenclatura exacta puede evolucionar con la implementación, pero la separación de responsabilidades es estable: el agente planificador define el análisis, el agente generador produce código, la capa de ejecución produce resultados verificables y el agente interpretador convierte esos resultados en lenguaje natural.

\begin{table}[htbp]
\centering
\small
\begin{tabularx}{\textwidth}{L{0.19\textwidth}L{0.27\textwidth}L{0.27\textwidth}Y}
\toprule
\textbf{Fase} & \textbf{Recibe} & \textbf{Genera} & \textbf{Utilidad en la memoria} \\
\midrule
Entrada normalizada & Consulta, tickers, rango temporal y rutas CSV. & \texttt{FinancialQueryInput} o estructura equivalente. & Fija qué datos puede usar el sistema y evita ambigüedades en la ejecución. \\
Agente planificador & Entrada normalizada y reglas de niveles A/B/C. & Plan de análisis con objetivo, métricas, salidas solicitadas y restricciones. & Permite evaluar si el modelo entendió la consulta antes de calcular nada. \\
Agente generador & Consulta original, entrada normalizada y plan analítico. & Script Python autocontenido ajustado al contrato. & Permite adaptar el cálculo a cada consulta sin perder validación. \\
Capa de ejecución & Código validado y datos históricos preparados. & Métricas, tablas y datos estructurados derivados de los CSV. & Separa cálculo reproducible de redacción en lenguaje natural. \\
Agente interpretador & Consulta original, plan y resultados estructurados. & Respuesta final comprensible, prudente y sin recomendación de inversión. & Permite valorar claridad, fidelidad a los datos y reconocimiento de limitaciones. \\
Trazabilidad & Artefactos anteriores y estado de ejecución. & Registro revisable de entrada, plan, cálculo, errores y respuesta. & Facilita auditoría, depuración y comparación entre modelos. \\
\bottomrule
\end{tabularx}
\caption{Material generado por el flujo y papel de cada agente.}
\label{tab:artefactos_flujo_agentes}
\end{table}

Esta lectura también ayuda a distinguir material de trabajo de resultados finales. Las tablas procesadas y las notas auxiliares se conservan como evidencias para reconstruir decisiones de diseño, parámetros de descarga, estructura de los CSV, criterios de evaluación y salidas intermedias. No se presentan como conclusiones por sí mismas, sino como soporte metodológico para que el flujo sea reproducible y revisable.
